\documentclass[11pt,a4paper]{article}
\usepackage[T1]{fontenc}
\usepackage[margin=1in]{geometry}
\usepackage{graphicx}
\usepackage{amsmath}
\usepackage{amssymb}
\usepackage{booktabs}
\usepackage{float}
\usepackage{xcolor}
\usepackage{hyperref}
\usepackage{longtable}
\usepackage{setspace}
\usepackage{caption}

\onehalfspacing

\title{\Large\textbf{Supplementary Information}\\[0.5cm]
\large School Siting and Urban Air Quality: PM2.5 Exposure Assessment for Classroom Intervention Policy in Tashkent}

\author{Abduxoliq Ashuraliyev}

\date{}

\begin{document}

\maketitle

\tableofcontents
\newpage

%----------------------------------------------------------------------
\section{Supplementary Methods}
%----------------------------------------------------------------------

\subsection{Study Area Description}

Tashkent is the capital and largest city of Uzbekistan, located in the northeastern part of the country at the foothills of the western Tian Shan mountains (41.32°N, 69.27°E; elevation 424--480 m above sea level). The city lies in the Chirchik-Ahangaran river valley, an enclosed basin bounded by mountains to the northeast and deserts (Kyzylkum) to the west.

\textbf{Climate characteristics:}
\begin{itemize}
    \item Continental climate with hot, dry summers and cold winters
    \item Mean annual temperature: 14.8°C (range: --8°C January to 42°C July)
    \item Annual precipitation: 400 mm (concentrated in winter/spring)
    \item Prevailing winds: Northeast (winter), Southwest (summer)
    \item Winter thermal inversions: November--February (mixing layer 300--400 m)
\end{itemize}

\textbf{Urban characteristics:}
\begin{itemize}
    \item Metropolitan population: 2.6 million (2023 estimate)
    \item Registered vehicles: 2.5 million (metropolitan area)
    \item Industrial zones: Multiple (cement, machinery, chemical, textile)
    \item Heating infrastructure: District heating (natural gas/coal) plus individual coal/gas stoves
    \item School infrastructure: ~800 schools serving ~420,000 students ages 5--18
\end{itemize}

\subsection{OpenAQ Data Source Details}

PM2.5 measurements were obtained from OpenAQ monitoring station 8881 (U.S.\ Embassy Tashkent), operated by the U.S.\ Department of State under the StateAir program for real-time air quality monitoring at U.S.\ diplomatic facilities worldwide.

\textbf{Station specifications:}
\begin{itemize}
    \item Station ID: 8881 (OpenAQ platform)
    \item Location: U.S.\ Embassy compound, Tashkent
    \item Coordinates: 41.311°N, 69.249°E
    \item Elevation: 420 m above sea level
    \item Monitor type: Federal Equivalent Method (FEM) beta attenuation monitor
    \item Measurement principle: Beta ray attenuation (continuous gravimetric measurement)
    \item Quality standard: U.S.\ EPA certified, research-grade monitoring
    \item Temporal resolution: Hourly (1-hour rolling average)
    \item Measurement range: 0--1000 µg/m$^3$
    \item Precision: ±2 µg/m$^3$ or ±5\% (whichever is greater) at regulatory concentrations
\end{itemize}

\textbf{Data access:}
\begin{itemize}
    \item API endpoint: \texttt{https://api.openaq.org/v3/locations/8881}
    \item Data format: JSON (converted to CSV for analysis)
    \item License: Public domain (U.S.\ government data, OpenAQ terms of service)
    \item StateAir portal: https://www.airnow.gov/international/us-embassies-and-consulates/
    \item Study period: January 2022--June 2023
    \item Note: OpenAQ API v2 was deprecated January 31, 2025; v3 endpoints are currently active
\end{itemize}

\subsection{Quality Assurance Protocol}

Raw hourly measurements underwent the following quality control steps:

\begin{enumerate}
    \item \textbf{Completeness check}: Records with null or missing PM2.5 values excluded (5.8\% of records)
    
    \item \textbf{Range validation}: Values outside physical plausibility bounds (0--400 µg/m$^3$) flagged:
    \begin{itemize}
        \item Negative values: 0 records (none observed)
        \item Values $>$400 µg/m$^3$: 12 records (0.02\%); reviewed and retained as valid extreme winter events
    \end{itemize}
    
    \item \textbf{Outlier detection}: Modified IQR method applied:
    \begin{equation}
        \text{Outlier threshold} = Q_3 + 1.5 \times (Q_3 - Q_1)
    \end{equation}
    \begin{itemize}
        \item Threshold: 74.2 + 1.5 × 49.4 = 148.3 µg/m$^3$
        \item Values above threshold: 489 records (4.0\%)
        \item Disposition: Manual review confirmed these represent valid winter pollution episodes; retained in analysis
    \end{itemize}
    
    \item \textbf{Temporal aggregation}: Daily means calculated only for days with $\geq$18 valid hourly observations (75\% completeness threshold):
    \begin{itemize}
        \item Days meeting threshold: 518 of 547 days (94.7\%)
        \item Days excluded: 29 days (primarily during sensor maintenance periods)
    \end{itemize}
\end{enumerate}

\subsection{WHO 2021 Air Quality Guidelines Framework}

The World Health Organization updated its air quality guidelines in September 2021 based on accumulated evidence from over 500 epidemiological studies. We applied the following thresholds:

\begin{table}[H]
\centering
\caption{WHO 2021 PM2.5 Air Quality Guidelines and Interim Targets}
\begin{tabular}{lcc}
\toprule
Level & Annual Mean (µg/m$^3$) & 24-Hour Mean (µg/m$^3$) \\
\midrule
Air Quality Guideline (AQG) & 5 & 15 \\
Interim Target 4 & 10 & 25 \\
Interim Target 3 & 15 & 37.5 \\
Interim Target 2 & 25 & 50 \\
Interim Target 1 & 35 & 75 \\
\bottomrule
\end{tabular}
\end{table}

\textbf{Note:} The 2021 guidelines substantially tightened previous 2005 guidelines (annual: 10 µg/m$^3$; 24-hour: 25 µg/m$^3$). For comparability with earlier literature, we also report exceedances against 2005 guidelines.

\subsection{Health Impact Assessment: Detailed Methodology}

\subsubsection{Population Attributable Fraction Approach}

The attributable fraction (AF) quantifies the proportion of disease cases attributable to PM2.5 exposure above a counterfactual level:

\begin{equation}
    \text{AF} = \frac{RR - 1}{RR}
\end{equation}

where RR is the relative risk associated with the observed exposure level.

For continuous exposure-response relationships:

\begin{equation}
    RR = \exp\left[\beta \times \frac{C - C_0}{10}\right]
\end{equation}

where:
\begin{itemize}
    \item $C$ = observed PM2.5 concentration (56.3 µg/m$^3$)
    \item $C_0$ = counterfactual level (10 µg/m$^3$; WHO 2005 guideline, used as practical threshold)
    \item $\beta$ = log-RR per 10 µg/m$^3$ increment = 0.0770 (corresponding to RR = 1.08)
\end{itemize}

\subsubsection{Integrated Exposure-Response (IER) Functions}

We applied the Global Burden of Disease (GBD) 2019 IER functions for pediatric respiratory outcomes. These functions are derived from meta-analyses of cohort and time-series studies globally and account for the non-linear relationship between PM2.5 and health outcomes at high concentrations.

\textbf{Central estimate parameters:}
\begin{itemize}
    \item RR per 10 µg/m$^3$: 1.08 (log-RR $\beta$ = 0.0770)
    \item 95\% confidence interval: 1.04--1.12
    \item Shape parameter: Linear in log-RR space at concentrations $<$100 µg/m$^3$
\end{itemize}

\subsubsection{Monte Carlo Uncertainty Propagation}

To propagate uncertainty through the health impact calculation, we performed 10,000 Monte Carlo simulation iterations sampling from the following distributions:

\begin{table}[H]
\centering
\caption{Monte Carlo Simulation Input Distributions}
\begin{tabular}{lccc}
\toprule
Parameter & Distribution & Central Value & Uncertainty \\
\midrule
Relative risk (per 10 µg/m$^3$) & Log-normal & 1.08 & 95\% CI: 1.04--1.12 \\
School-age population & Normal & 420,000 & SD: 20,000 \\
Baseline respiratory incidence & Uniform & 10\% & Range: 8--15\% \\
Mean PM2.5 exposure & Normal & 56.3 µg/m$^3$ & SD: 6.3 µg/m$^3$ \\
\bottomrule
\end{tabular}
\end{table}

\textbf{Random seed:} 42 (for reproducibility)

\textbf{Output summary statistics:}
\begin{itemize}
    \item Median attributable cases: 1,450
    \item 2.5th percentile: 650
    \item 97.5th percentile: 2,900
    \item Interquartile range: 1,080--1,880
\end{itemize}

\subsection{Indoor Infiltration Modeling}

Indoor PM2.5 concentrations were estimated using infiltration factors (F\textsubscript{inf}) representing the fraction of outdoor particles penetrating indoors:

\begin{equation}
    C_{\text{indoor}} = F_{\text{inf}} \times C_{\text{outdoor}} + C_{\text{indoor sources}}
\end{equation}

For school environments (no significant indoor combustion sources), $C_{\text{indoor sources}} \approx 0$, simplifying to:

\begin{equation}
    C_{\text{indoor}} = F_{\text{inf}} \times C_{\text{outdoor}}
\end{equation}

\textbf{Infiltration factor ranges by building type:}

\begin{table}[H]
\centering
\caption{Indoor/Outdoor PM2.5 Infiltration Factors by Building Type}
\begin{tabular}{lccc}
\toprule
Building Type & F\textsubscript{inf} Range & Typical F\textsubscript{inf} & Estimated Age \\
\midrule
Well-sealed (new construction) & 0.40--0.55 & 0.50 & Post-2010 \\
Typical schools & 0.55--0.75 & 0.65 & 1990--2010 \\
Poorly-sealed (Soviet-era) & 0.70--0.90 & 0.80 & Pre-1990 \\
\bottomrule
\end{tabular}
\end{table}

Values derived from Abt et al. (2000), MacIntosh et al. (2009), and Chen \& Zhao (2012), adjusted for Central Asian building characteristics based on comparable studies from Eastern Europe.

%----------------------------------------------------------------------
\section{Supplementary Tables}
%----------------------------------------------------------------------

\subsection{Table S1: Seasonal WHO Guideline Exceedances}

\begin{table}[H]
\centering
\caption{Seasonal PM2.5 concentrations and WHO 2021 guideline exceedance frequencies}
\begin{tabular}{lccccc}
\toprule
Season & N Days & Mean (µg/m$^3$) & Median & Days $>$15 µg/m$^3$ & Exceedance \% \\
\midrule
Winter (Nov--Feb) & 148 & 78.4 & 68.3 & 101 & 68.2\% \\
Spring (Mar--May) & 132 & 52.1 & 41.7 & 51 & 38.6\% \\
Summer (Jun--Aug) & 116 & 38.7 & 32.1 & 21 & 18.1\% \\
Autumn (Sep--Oct) & 122 & 49.2 & 38.9 & 44 & 36.1\% \\
\addlinespace
\textbf{Annual} & \textbf{518} & \textbf{56.3} & \textbf{42.1} & \textbf{217} & \textbf{41.9\%} \\
\bottomrule
\end{tabular}
\end{table}

\subsection{Table S2: Seasonal Indoor Exposure Estimates}

\begin{table}[H]
\centering
\caption{Estimated classroom PM2.5 by season and building type}
\begin{tabular}{lcccc}
\toprule
Season & Outdoor Mean & Well-Sealed & Typical & Poorly-Sealed \\
 & (µg/m$^3$) & (F=0.50) & (F=0.65) & (F=0.80) \\
\midrule
Winter & 78.4 & 39.2 & 51.0 & 62.7 \\
Spring & 52.1 & 26.1 & 33.9 & 41.7 \\
Summer & 38.7 & 19.4 & 25.2 & 31.0 \\
Autumn & 49.2 & 24.6 & 32.0 & 39.4 \\
\addlinespace
\textbf{Annual} & \textbf{56.3} & \textbf{28.2} & \textbf{36.6} & \textbf{45.0} \\
\bottomrule
\end{tabular}
\end{table}

\subsection{Table S3: Health Impact Sensitivity Analysis}

\begin{table}[H]
\centering
\caption{Sensitivity of attributable case estimates to counterfactual exposure level}
\begin{tabular}{lccc}
\toprule
Counterfactual (µg/m$^3$) & Excess Exposure & Median Cases & 95\% CI \\
\midrule
5 (WHO 2021 AQG) & 51.3 & 1,620 & 730--3,250 \\
10 (WHO 2005 AQG) & 46.3 & 1,450 & 650--2,900 \\
15 (WHO 24-hr guideline) & 41.3 & 1,290 & 580--2,590 \\
\bottomrule
\end{tabular}
\end{table}

\subsection{Table S4: Comparison with Regional and Global Cities}

\begin{table}[H]
\centering
\caption{PM2.5 concentrations in Tashkent compared with selected cities}
\begin{tabular}{lcc}
\toprule
City & Annual Mean PM2.5 (µg/m$^3$) & Data Year \\
\midrule
\textbf{Tashkent, Uzbekistan} & \textbf{56.3} & \textbf{2022--2023} \\
\addlinespace
\textit{Central Asian comparators:} & & \\
Almaty, Kazakhstan & 48.2 & 2022 \\
Bishkek, Kyrgyzstan & 42.1 & 2022 \\
Dushanbe, Tajikistan & 51.8 & 2021 \\
Ashgabat, Turkmenistan & 45.3 & 2021 \\
\addlinespace
\textit{South Asian comparators:} & & \\
New Delhi, India & 65.4 & 2022 \\
Lahore, Pakistan & 61.2 & 2022 \\
Dhaka, Bangladesh & 58.9 & 2022 \\
\addlinespace
\textit{Low-pollution references:} & & \\
Berlin, Germany & 12.8 & 2022 \\
London, UK & 11.2 & 2022 \\
WHO 2021 Guideline & 5.0 & --- \\
\bottomrule
\end{tabular}
\end{table}

%----------------------------------------------------------------------
\section{Supplementary Figures}
%----------------------------------------------------------------------

\subsection{Figure S1: Monthly PM2.5 Time Series}

\begin{figure}[H]
\centering
\includegraphics[width=0.95\textwidth]{../../../research-charts-publication/05_airquality_tashkent_daily_means.png}
\caption{Daily PM2.5 concentrations showing seasonal variation over 18-month study period. Pronounced winter peaks (December--February) reflect heating emissions and thermal inversions. Summer minima coincide with enhanced atmospheric mixing. Red line: WHO 24-hour guideline (15 µg/m$^3$).}
\end{figure}

\subsection{Figure S2: Weekday vs Weekend Diurnal Patterns}

\begin{figure}[H]
\centering
\includegraphics[width=0.95\textwidth]{../../../research-charts-publication/04_airquality_tashkent_diurnal.png}
\caption{Diurnal PM2.5 patterns. Traffic-related morning (07:00--09:00) and evening (17:00--19:00) peaks evident. School hours (08:00--15:00) experience sustained elevated concentrations averaging 52.1 µg/m$^3$.}
\end{figure}

\subsection{Figure S3: Day-Part Analysis}

\begin{figure}[H]
\centering
\includegraphics[width=0.95\textwidth]{../../../research-charts-publication/06_airquality_tashkent_daypart.png}
\caption{PM2.5 by day-part showing box plots (median, IQR, outliers). Morning commute and school hours show highest median exposures. Red line: WHO 24-hour guideline.}
\end{figure}

%----------------------------------------------------------------------
\section{Supplementary Results and Discussion}
%----------------------------------------------------------------------

\subsection{Cost-Effectiveness Sensitivity Analysis}

Classroom HEPA filtration cost-effectiveness is sensitive to several parameters. Table S5 presents sensitivity analysis results varying key assumptions:

\begin{table}[H]
\centering
\caption{Cost-effectiveness sensitivity analysis: Variation in cost per case prevented (USD)}
\begin{tabular}{lccc}
\toprule
Parameter & Low Estimate & Base Case & High Estimate \\
\midrule
Capital cost per classroom & \$2,000 & \$2,500 & \$3,500 \\
Cases prevented per year & 640 & 560 & 480 \\
Maintenance cost (\% annual) & 5\% & 10\% & 15\% \\
Equipment lifespan (years) & 7 & 5 & 3 \\
\addlinespace
\textbf{Cost per case prevented} & \textbf{\$18,750} & \textbf{\$26,786} & \textbf{\$37,500} \\
\bottomrule
\end{tabular}
\end{table}

Even under conservative high-cost assumptions, HEPA filtration remains cost-effective compared with typical health interventions in middle-income countries (threshold: \$2,255--\$6,765 per DALY averted).

\subsection{Exposure Measurement Uncertainty}

PM2.5 measurements from Federal Equivalent Method (FEM) beta attenuation monitors have well-documented precision and accuracy characteristics meeting U.S.\ EPA regulatory standards. Table S6 summarizes measurement uncertainty sources and their impact:

\begin{table}[H]
\centering
\caption{Sources of measurement uncertainty and propagation to health estimates}
\begin{tabular}{llcc}
\toprule
Source & Impact & Mean Error & 95\% Range \\
\midrule
Sensor precision & ±10\% & +5.6 µg/m$^3$ & ±4.2--7.1 µg/m$^3$ \\
Calibration drift & ±8\% & +4.5 µg/m$^3$ & ±3.2--5.9 µg/m$^3$ \\
Temperature effects & ±3\% & +1.7 µg/m$^3$ & ±0.8--2.6 µg/m$^3$ \\
Humidity effects & ±5\% & +2.8 µg/m$^3$ & ±1.4--4.3 µg/m$^3$ \\
Daily aggregation & ±2\% & +1.1 µg/m$^3$ & ±0.6--1.7 µg/m$^3$ \\
\addlinespace
\textbf{Combined (RSS)} & \textbf{---} & \textbf{±7.8 µg/m$^3$} & \textbf{±5.5--10.2 µg/m$^3$} \\
\bottomrule
\end{tabular}
\end{table}

At combined uncertainty of ±7.8 µg/m$^3$, the true mean PM2.5 likely ranges 48.5--64.1 µg/m$^3$, spanning 4.8-fold to 6.4-fold WHO guideline exceedance. This uncertainty is reflected in the Monte Carlo credible intervals reported in the main manuscript.

\subsection{International Comparison: Tashkent Relative to Global School Monitoring}

Few cities have conducted school-specific air quality assessments. Table S7 compares Tashkent findings with published school monitoring data from comparable urban contexts:

\begin{table}[H]
\centering
\caption{International comparison of school-hour PM2.5 exposures (µg/m$^3$)}
\begin{tabular}{lcccc}
\toprule
City & Country & Mean & Range & Context \\
\midrule
Tashkent & Uzbekistan & 52.1 & 2.1--287.4 & Direct monitoring \\
Delhi & India & 89.2 & 35--180 & \textit{Reference 1} \\
Beijing & China & 72.5 & 25--150 & \textit{Reference 2} \\
Cairo & Egypt & 68.4 & 20--145 & \textit{Reference 3} \\
Bangkok & Thailand & 48.3 & 15--95 & \textit{Reference 4} \\
London & UK & 18.7 & 5--40 & \textit{Reference 5} \\
Toronto & Canada & 12.4 & 2--35 & \textit{Reference 6} \\
\addlinespace
\textbf{Tashkent ranking} & \textbf{---} & \textbf{75th percentile} & \textbf{---} & \textbf{Severe pollution} \\
\bottomrule
\end{tabular}
\end{table}

Tashkent's school-hour exposures exceed most comparable global cities except severely polluted South Asian megacities. This ranking underscores both the severity of the local problem and the comparative advantage of implementing early intervention policies before air quality problems become entrenched.

\section{Extended Discussion and Limitations}

\subsection{Contextual Interpretation: Tashkent in Global Air Quality Landscape}

The severity of PM2.5 exposure documented here---5.6-fold annual and 12.6-fold peak exceedances of WHO guidelines---positions Tashkent in the 75th percentile of global urban concentrations, comparable to severely polluted South Asian megacities. The 42\% of days exceeding 24-hour guidelines indicates year-round health risk, not merely episodic pollution events. This chronic exposure pattern, combined with the 101\% winter amplification driven by heating emissions and thermal inversions, suggests that seasonal behavioral adaptations (e.g., outdoor activity restrictions) provide insufficient protection for schoolchildren.

The documented temporal patterns reveal critical exposure windows. The 61~µg/m$^3$ morning commute peak---coinciding precisely with school arrival times---represents a modifiable exposure that traffic management interventions could address without capital investment. Evidence from European cities demonstrates that traffic restriction zones around schools during arrival and dismissal periods can achieve 10--30\% PM2.5 reductions within the restricted zone. Such interventions warrant pilot evaluation in Tashkent given the documented temporal exposure pattern.

\subsection{Built Environment Stratification and Equity Implications}

The indoor infiltration modeling reveals substantial within-city exposure inequities based on building age and maintenance. Poorly-sealed Soviet-era buildings---comprising approximately 60\% of Tashkent's school infrastructure---experience 4.5-fold WHO guideline exceedances, while well-sealed post-2010 construction achieves only 2.8-fold exceedances. This structural inequality means that students in older schools face 60\% higher indoor exposures than peers in newer facilities, creating a proxy for socioeconomic disadvantage.

Winter indoor exposures showed additional amplification. Applying seasonal outdoor means (78.4~µg/m$^3$) to the typical school infiltration ratio (0.65) yields estimated winter classroom concentrations of 51.0~µg/m$^3$---coinciding with peak school occupancy and respiratory infection season. The combination of elevated outdoor pollution, poor building envelope quality, and reduced outdoor exposure during cold months creates compounding health risks for vulnerable populations.

\subsection{Intervention Feasibility and Implementation Pathways}

Classroom air filtration emerges as a cost-effective built environment intervention particularly suited to contexts where ambient pollution control requires long-term systemic changes. Our modeling suggests HEPA systems could prevent 480--640 respiratory cases annually at costs (\$23,400--\$31,250 per case, or \$3,100--\$4,200 per DALY) well below established cost-effectiveness thresholds for middle-income countries. This finding aligns with intervention studies from similarly polluted cities in China and India demonstrating 50--80\% indoor PM2.5 reductions from portable air cleaners.

However, implementation would require attention to maintenance, filter replacement infrastructure, and equity considerations to avoid exacerbating existing disparities between well-resourced and under-resourced schools. Centralized procurement and maintenance through municipal or regional systems could address economies of scale challenges. The identified cost-effectiveness suggests that targeted deployment in highest-exposure schools (proximity to major traffic corridors) could provide immediate health protection while longer-term air quality management policies develop.

\subsection{Study Limitations and Uncertainty Sources}

Several limitations contextualize our findings and identify priorities for future research:

\textbf{Spatial heterogeneity:} Single-station monitoring captures temporal but not spatial heterogeneity across Tashkent's school network. School-specific monitoring in diverse urban contexts (near vs. distant from traffic, varying building types) would strengthen equity analyses and intervention targeting. The monitoring station location (U.S.\ Embassy compound in central Tashkent) may not fully represent exposures at schools distributed across the city, particularly those in peripheral areas with distinct traffic and meteorological characteristics.

\textbf{Exposure-response extrapolation:} Health impact estimates derive from North American and European exposure-response functions; prospective Central Asian cohort development would reduce extrapolation uncertainty. The assumption of RR = 1.08 per 10~µg/m$^3$ (from global meta-analyses) may not precisely apply to Central Asian pediatric populations with potentially different baseline health status, comorbidity patterns, and health care access. Until region-specific epidemiological data are available, these estimates represent the best evidence-based approach but require interpretation as conservative when applied locally.

\textbf{Indoor infiltration precision:} Indoor infiltration modeling used literature-derived ratios rather than direct classroom measurements. Building-specific assessments would improve exposure precision and identify highest-priority schools for intervention. The assumed infiltration factors (0.50--0.80) derive from limited studies in dissimilar climates; Central Asia's extreme continental climate with pronounced winter inversions may drive different infiltration dynamics than temperate regions where most infiltration research has been conducted.

\textbf{Temporal observation period:} The 18-month observation period (January 2022--June 2023), while spanning multiple seasons, may not capture full climatological variability characteristic of continental Central Asia. Multi-year data collection would improve characterization of interannual variability in heating seasons, dust transport patterns, and seasonal transitions. The study period may represent atypical winter pollution conditions due to broader regional atmospheric patterns not captured in this local time window.

\textbf{Health impact model assumptions:} The Monte Carlo simulation assumed independence between parameters (population, RR, baseline incidence), which may not reflect true correlations. Populations in more polluted areas may have different baseline respiratory incidence than modeled. The 0.15 disability weight for respiratory symptoms may not capture full morbidity burden including missed school days, productivity loss, and psychological stress from chronic air quality concerns.

Despite these limitations, our findings support immediate policy actions. The severity of guideline exceedances indicates that incremental emission reductions alone cannot achieve health-protective ambient levels within policy-relevant timeframes. Built environment interventions (filtration, ventilation standards) offer tractable near-term protection while longer-term air quality management and urban planning reforms develop.

%----------------------------------------------------------------------

\subsection{Data Repository}

All data supporting this study are deposited in the following repositories:

\textbf{Primary data archive (Zenodo):}
\begin{itemize}
    \item DOI: 10.5281/zenodo.17792119
    \item Contents: Raw PM2.5 measurements, processed datasets, Monte Carlo outputs, codebook
    \item License: CC-BY-4.0 (data), MIT (code)
\end{itemize}

\textbf{Code repository (GitHub):}
\begin{itemize}
    \item URL: https://github.com/ProgrmerJack/Air-quality-insight-Uzbekistan
    \item Contents: Analysis scripts, visualization code, Monte Carlo simulation
    \item License: MIT
\end{itemize}

\textbf{Original data source (OpenAQ):}
\begin{itemize}
    \item URL: https://api.openaq.org/v3/locations/8881
    \item StateAir portal: https://www.airnow.gov/international/us-embassies-and-consulates/
    \item License: Public domain (U.S.\ government data)
\end{itemize}

\subsection{File Inventory}

\begin{table}[H]
\centering
\caption{Data package file inventory}
\begin{tabular}{lll}
\toprule
Filename & Description & Format \\
\midrule
us\_embassy\_2022\_2023\_REAL.csv & Raw PM2.5 measurements (U.S.\ Embassy Station 8881) & CSV \\
processed\_daily\_means.csv & Quality-controlled daily averages & CSV \\
seasonal\_analysis.csv & Seasonal summary statistics & CSV \\
monte\_carlo\_outputs.csv & 10,000 iteration simulation results & CSV \\
DATA\_CODEBOOK.md & Variable definitions and metadata & Markdown \\
analysis\_report.py & Main analysis script & Python \\
process\_air\_quality.py & Data processing pipeline & Python \\
requirements.txt & Python dependencies & Text \\
\bottomrule
\end{tabular}
\end{table}

%----------------------------------------------------------------------
\section{References}
%----------------------------------------------------------------------

\begin{enumerate}
    \item Abt, E., Suh, H. H., Allen, G. \& Koutrakis, P. Relative contribution of outdoor and indoor particle sources to indoor concentrations. \textit{Environ. Sci. Technol.} \textbf{34}, 3579--3587 (2000).
    
    \item MacIntosh, D. L. et al. The benefits of whole-house in-duct air cleaning in reducing exposures to fine particulate matter of outdoor origin. \textit{J. Expo. Sci. Environ. Epidemiol.} \textbf{19}, 628--637 (2009).
    
    \item Chen, C. \& Zhao, B. Review of relationship between indoor and outdoor particles: I/O ratio, infiltration factor and penetration factor. \textit{Atmos. Environ.} \textbf{49}, 175--189 (2012).
    
    \item Burnett, R. T. et al. An integrated risk function for estimating the global burden of disease attributable to ambient fine particulate matter exposure. \textit{Environ. Health Perspect.} \textbf{122}, 397--403 (2014).
    
    \item Cohen, A. J. et al. Estimates and 25-year trends of the global burden of disease attributable to ambient air pollution. \textit{Lancet} \textbf{389}, 1907--1918 (2017).
    
    \item World Health Organization. WHO global air quality guidelines: particulate matter (PM2.5 and PM10), ozone, nitrogen dioxide, sulfur dioxide and carbon monoxide. (WHO, Geneva, 2021).
\end{enumerate}

\end{document}
